\section{Computational study}\label{sec:experiments}
\subsection{Protocol, computation, and economic comparisons}
The recorded protocol fixes $T=8$, $\beta=0.97$, $z\in\{0,1,2\}$, and $D_z=z+D_0$, where $D_0\in\{0,1,2,3\}$ has probabilities $(0.1,0.3,0.4,0.2)$. The transition is
\[
 P=\begin{pmatrix}0.70&0.25&0.05\\0.15&0.70&0.15\\0.05&0.25&0.70\end{pmatrix}.
\]
The coefficients are $c=0.3,h=0.6,p=0.65,\kappa=1.2,m=1.5,\lambda=0.6$, and $r_z=0.6+1.4z$. Stock is $0,0.5,\ldots,6$, contract is $0,0.1,\ldots,1$, and the initial state is $(1,0.5)$. Demand translation makes $\Delta L\leq\Delta z$, while $\Delta r=1.4\Delta z>1.2\Delta z$. The transition rows are stochastically ordered. Thus the structural theorem applies to exactly the model computed.

All five methods use unrestricted state-aware stock tables. Joint contract and expanded-action MDP enumerate all 143 action pairs. Frozen contract fixes the economic contract at $0.5$ and has 13 distinct stock actions. The capacity-matched gate enumerates 11 dummy labels for each stock choice but leaves the economic contract and next inherited state at $0.5$; it still has only 13 distinct economic actions. The physical oracle minimizes physical cost. Under the recorded greatest-pair convention it selects contract one on contract ties; only its physical-cost comparison is invariant to this implementation convention.

Table~\ref{tab:exact} reports exact discounted expectations. Joint and expanded-action values and policies agree at every nonterminal state, as do frozen and dummy-gate values and policies. The independent Fraction implementation checks 1,320 state--method records; the monotonicity check covers 416 adjacent state pairs. The hull critic agrees at all 264 states. Exact Bellman, terminal, greedy, and accounting defects are zero for the finite model.

Net reward improves from $-5.947058$ to $-5.002311$, a gain of $0.944747$. Physical cost falls by $0.213204$, from $9.587378$ to $9.374174$. Nevertheless, the fill rate falls from $0.925926$ to $0.882109$. The physical oracle costs $8.506095$, below both contract policies. These results have a direct economic reading: the joint contract changes the cash-flow value of carrying stock, and it economizes on physical cost partly by accepting more lost demand. Improved welfare is not a synonym for higher service.

The full accounting decomposition appears in Table~\ref{tab:ledger}. For example, the joint policy earns more premium but also pays more liability, maintenance, and adjustment. Equation~\eqref{eq:identity} accounts for all these changes. No historical total-cost field is relabeled as theta-free physical cost.

The simulation uses 64 batches of 4,096 episodes with common exogenous paths across methods and the recorded seed sequence. Two-sided paired Student intervals use 63 degrees of freedom and quantify simulation error only. The paired joint-minus-frozen net-reward difference is $0.943729$ with interval $[0.939821,0.947636]$; the physical-cost difference is $-0.212112$ with interval $[-0.214933,-0.209291]$. Exact expectations, rather than Monte Carlo significance, establish the finite-model ordering. Batch-level outputs also retain undiscounted service counts.

\subsection{Contract refinement, conditioning, and diffusion certificates}
Table~\ref{tab:contractgrid} reports the recorded nested contract grids with 10, 20, 40, 80, and 160 intervals. Each row solves the same stock-grid model exactly. Theorem~\ref{thm:quadratic} provides the continuous-contract enclosure, with error bounds decreasing from $0.047158$ to $0.00018421$. These are analytic error budgets, not thresholds fitted to observed grid differences. A first-order Lipschitz budget is retained in the machine-readable output for comparison.

The conditioning experiment uses 64 batches of 2,048 identical observations for each parameterization. For $u$ uniform on $[-1,1]$ and standard normal $\xi$, the observed payoff is $(u+\sqrt{2\varepsilon}\xi)^2/2$. We fit $ku^2/2+b$ and $\zeta u^2/(2\sqrt{2\varepsilon})+b$ by unregularized least squares. The resulting shadow prices are $ku$ and $\zeta u/\sqrt{2\varepsilon}$. Table~\ref{tab:conditioning} shows matching mean-square errors, with maximum slope disagreement below $1.2\times10^{-15}$. The volatility design becomes less well-conditioned at the smallest noise level, but the exact least-squares fit does not display a shadow-price accuracy advantage for the direct parameterization. No artificial target-noise floor is added to create one.

We additionally compile the smooth critics of Corollary~\ref{cor:neural} at $\tau=2^{-m}$ for $m=4,8,12,16,20$. They use at most 27 cubic-ReLU units per state--time critic. Exact rational evaluation checks every one of the 264 nonterminal grid states at each level. All five induced policies coincide with the exact grid optimizer at these states, giving zero observed grid-policy loss. The proved uniform policy budget, which also covers initial inherited contracts between grid points, decreases from $0.0776069$ to $1.18419\times10^{-6}$. EC.4 reports the weights, finite-model interpretation, and complete audit locations. This is constructive approximation and validation, not a newly trained neural solver.

Table~\ref{tab:pde} instantiates \eqref{eq:analytic} at the five recorded perturbation levels. It reports separate terminal and lateral defects and the absolute policy-loss bound. The bound decreases to approximately $4.38690\times10^{-6}$. The state--time--action study in Table~\ref{tab:mesh} is an additional deterministic scheme check, separate from the recorded analytic sequence. It refines 25 to 1,089 spatial nodes, 4 to 94 time steps, and 9 to 65 actions per coordinate. The observed all-node error decreases from $0.0445124$ to $0.00742613$, within the proved scheme-specific bounds. Floating-point observed errors are diagnostics; the bounds follow from the analytic truncation calculation and nonexpansive propagation.

\subsection{Evidence retained from earlier versions}
The complete earlier manuscript, proofs, figures, numerical outputs, release archive, and reproducibility checks remain available in the historical snapshot. EC.5 identifies their relationship to the present paper. Historical large-neural inventory residual ratios and greedy-gap summaries are not reclassified as uniform certificates: they lack the absolute parabolic-boundary and interior enclosures required by Theorem~\ref{thm:certificate}. The finite model, the manufactured diffusion, and those sampled neural diagnostics are different mathematical objects. A successful exact calculation on one does not supply missing hypotheses for another.

\section{Operational implications and reproducibility}\label{sec:conclusion}
A costly contract has operational content when its cash flows are priced externally and its persistence affects future decisions. The premium--liability condition explains when stress raises the chosen commitment rather than inducing the provider to lower an internal score. The envelope representation turns the inherited contract into an ordered threshold decision and supplies an approximation guarantee without global concavity. These results are fully compatible with ordinary expanded-action stochastic control.

The numerical evidence supports a specific, reproducible conclusion: joint contracting improves net reward and physical cost over the recorded fixed-contract policy, while sacrificing fill rate and remaining above the physical-cost lower bound. For continuous reconfiguration, direct value gradients address a conditioning issue, whereas residual, greedy, terminal, and lateral defects determine certification. These roles should be kept distinct in algorithm design and in operational interpretation.

\paragraph{Code and data availability.}
The revision package contains executable integer and independently written rational dynamic programs, all policy and accounting records, paired simulation and conditioning outputs, analytic and grid-refinement calculations, and a point-by-point response. The entry point is \texttt{revisions/or-r4-20260920/execute.py}; generated evidence is in its \texttt{results} directory. The frozen protocol is retained at \texttt{revisions/or-r2-20260920/protocol.json}. No empirical data are needed for the new experiments. Earlier data and source artifacts are preserved with their original interpretation in the historical snapshot, not silently pooled with the new results.

\section*{Acknowledgments}
Author-identifying acknowledgments are omitted for anonymous review.

\begin{thebibliography}{9}
\bibitem[Federgruen and Heching(1999)]{FedergruenHeching1999}
Federgruen A, Heching A (1999) Combined pricing and inventory control under uncertainty. \emph{Operations Research} 47(3):454--475.
\bibitem[Han et al.(2018)Han, Jentzen, and E]{HanJentzenE2018}
Han J, Jentzen A, E W (2018) Solving high-dimensional partial differential equations using deep learning. \emph{Proceedings of the National Academy of Sciences} 115(34):8505--8510.
\bibitem[Kushner and Dupuis(2001)]{KushnerDupuis2001}
Kushner HJ, Dupuis P (2001) \emph{Numerical Methods for Stochastic Control Problems in Continuous Time}, 2nd ed. (Springer, New York).
\bibitem[Ma et al.(1994)Ma, Protter, and Yong]{MaProtterYong1994}
Ma J, Protter P, Yong J (1994) Solving forward-backward stochastic differential equations explicitly---a four step scheme. \emph{Probability Theory and Related Fields} 98:339--359.
\bibitem[Topkis(1978)]{Topkis1978}
Topkis DM (1978) Minimizing a submodular function on a lattice. \emph{Operations Research} 26(2):305--321.
\end{thebibliography}
\clearpage
\begin{table}[!ht]
\centering
\caption{Exact discounted operating outcomes}
\label{tab:exact}
\begin{tabular}{lrrrr}
\toprule
Method & Net reward & Physical cost & Fill rate & Actions\\
\midrule
Joint contract & -5.002311 & 9.374174 & 0.882109 & 143\\
Expanded-action MDP & -5.002311 & 9.374174 & 0.882109 & 143\\
Frozen contract & -5.947058 & 9.587378 & 0.925926 & 13\\
Dummy gate & -5.947058 & 9.587378 & 0.925926 & 13\\
Physical oracle & -11.972030 & 8.506095 & 0.703704 & 143\\
\bottomrule
\end{tabular}
\par\smallskip
\begin{minipage}{0.98\textwidth}\small All expectations are exact rational calculations; displayed decimals are rounded. Actions count distinct economic choices. The physical oracle minimizes physical cost, not the convention-dependent net reward shown here.\end{minipage}
\end{table}

\begin{table}[!ht]
\centering
\caption{Contract cash-flow decomposition}
\label{tab:ledger}
\begin{tabular}{lrrrr}
\toprule
Method & Premium & Liability & Maintenance & Adjustment\\
\midrule
Joint contract & 9.889658 & 1.001859 & 4.308149 & 0.207786\\
Expanded-action MDP & 9.889658 & 1.001859 & 4.308149 & 0.207786\\
Frozen contract & 7.208555 & 0.865027 & 2.703208 & 0.000000\\
Dummy gate & 7.208555 & 0.865027 & 2.703208 & 0.000000\\
Physical oracle & 14.417109 & 6.920212 & 10.812832 & 0.150000\\
\bottomrule
\end{tabular}
\par\smallskip
\begin{minipage}{0.98\textwidth}\small Net reward equals premium minus physical cost, liability, maintenance, and adjustment. The oracle uses contract one under the recorded tie convention.\end{minipage}
\end{table}

\clearpage
\begin{table}[!ht]
\centering
\caption{Exact grid values and continuous-contract enclosures}
\label{tab:contractgrid}
\begin{tabular}{rrrr}
\toprule
Intervals & Grid value & Loss bound & Value upper bound\\
\midrule
10 & -5.00231055 & 0.04715774 & -4.95515281\\
20 & -4.99648554 & 0.01178944 & -4.98469611\\
40 & -4.99331402 & 0.00294736 & -4.99036666\\
80 & -4.99288305 & 0.00073684 & -4.99214621\\
160 & -4.99278573 & 0.00018421 & -4.99260152\\
\bottomrule
\end{tabular}
\par\smallskip
\begin{minipage}{0.98\textwidth}\small Theorem~\ref{thm:quadratic} bounds contract discretization for the unchanged stock set. The lower endpoint is the grid value; the upper endpoint adds the proved bound.\end{minipage}
\end{table}

\begin{table}[!ht]
\centering
\caption{Identical-data least-squares conditioning experiment}
\label{tab:conditioning}
\begin{tabular}{rrrrr}
\toprule
$\varepsilon$ & Direct $P$ MSE &  $P$ MSE via $Z$ & Direct condition & $Z$ condition\\
\midrule
0.1 & 0.00131353 & 0.00131353 & 6.901 & 3.461\\
0.01 & 0.00010902 & 0.00010902 & 6.901 & 2.980\\
0.001 & 0.00001052 & 0.00001052 & 6.901 & 7.653\\
0.0001 & 0.00000104 & 0.00000104 & 6.901 & 23.719\\
\bottomrule
\end{tabular}
\par\smallskip
\begin{minipage}{0.98\textwidth}\small Means over 64 batches, each with 2,048 pairs. Integrated shadow-price MSE uses uniform test states on $[-1,1]$. Condition numbers are spectral design-matrix condition numbers.\end{minipage}
\end{table}

\clearpage
\begin{table}[!ht]
\centering
\caption{Absolute certificates for the stopped diffusion}
\label{tab:pde}
\begin{tabular}{rrrrrr}
\toprule
$m$ & $\delta_H$ bound & $\delta_T$ & $\delta_\partial$ & $\delta_A$ & Policy-loss bound\\
\midrule
4 & 0.147851563 & 0 & 0 & 0 & 0.295703125\\
8 & 0.009000397 & 0 & 0 & 0 & 0.018000793\\
12 & 0.000561586 & 0 & 0 & 0 & 0.001123172\\
16 & 0.000035095 & 0 & 0 & 0 & 0.000070191\\
20 & 0.000002193 & 0 & 0 & 0 & 0.000004387\\
\bottomrule
\end{tabular}
\par\smallskip
\begin{minipage}{0.98\textwidth}\small $e=2^{-m}$. The entire parabolic boundary is exact, and the actor is analytically greedy. These are manufactured analytic critics, not fitted neural inventory models.\end{minipage}
\end{table}

\begin{table}[!ht]
\centering
\caption{Joint space--time--action refinement}
\label{tab:mesh}
\begin{tabular}{rrrrr}
\toprule
Spatial nodes & Time steps & Actions per axis & Node error & Proved bound\\
\midrule
25 & 4 & 9 & 0.04451240 & 0.34765625\\
81 & 10 & 17 & 0.02567135 & 0.16503906\\
289 & 29 & 33 & 0.01404964 & 0.08032227\\
1089 & 94 & 65 & 0.00742613 & 0.03961182\\
\bottomrule
\end{tabular}
\par\smallskip
\begin{minipage}{0.98\textwidth}\small Observed errors are maximized over every computed state--time node. The bounds are supplied by Proposition~\ref{prop:scheme}; they do not certify interpolation or a continuous-state actor.\end{minipage}
\end{table}

\end{document}
